% Preámbulo común del manual del estudiante (main.tex) y de la guía del instructor
% (guia_instructor.tex). Se carga después de \documentclass y de marco-academic.

\usepackage{longtable}
\usepackage{calc}
\usepackage{float}
\usepackage{bookmark}
\usepackage{amssymb}
\usepackage{multirow}
\usepackage{ifthen}
\usepackage{tikz}
\usetikzlibrary{arrows.meta,positioning,calc,decorations.pathreplacing,decorations.pathmorphing,shapes.geometric,fit,backgrounds,babel}

\setcounter{secnumdepth}{2}
\setcounter{tocdepth}{1}
\RedeclareSectionCommand[tocnumwidth=2.5em]{section}
\setlength{\LTleft}{0pt}
% Con \LTright rígido, el \caption de longtable desborda exactamente \linewidth.
\setlength{\LTright}{0pt plus 1fill}
\setlength{\tabcolsep}{4pt}
\renewcommand{\arraystretch}{1.14}

% ---------------------------------------------------------------------------
% Numeración por laboratorio: 0.x (común), 1.x (Windows), 2.x (Ubuntu), A.x.
% Cada laboratorio reinicia secciones y pasos; el prefijo de paso (W/U) evita
% confusiones al citar «el paso W12» frente a «el paso U12».
% ---------------------------------------------------------------------------
\newcommand{\labnum}{0}
\renewcommand{\thesection}{\labnum.\arabic{section}}
\newcounter{paso}
\newcommand{\pasoprefijo}{P}
\renewcommand{\thepaso}{\pasoprefijo\arabic{paso}}
\newcommand{\IniciarLaboratorio}[2]{% #1 número (0,1,2,A), #2 prefijo de paso
  \renewcommand{\labnum}{#1}\setcounter{section}{0}%
  \renewcommand{\pasoprefijo}{#2}\setcounter{paso}{0}}

% Portadilla de laboratorio: aparece en el índice y en los marcadores del PDF.
\newcommand{\PortadillaLaboratorio}[3]{% #1 kicker, #2 título, #3 entradilla
  \clearpage
  \phantomsection
  \addcontentsline{toc}{part}{#2}%
  \begin{tcolorbox}[enhanced,frame hidden,colback=MAPrimaryDark,coltext=white,arc=2mm,
      left=7mm,right=7mm,top=6mm,bottom=6mm]
    {\sffamily\bfseries\footnotesize\color{white!80!MAPrimarySoft}\MakeUppercase{#1}}\par\vspace{2mm}
    {\sffamily\bfseries\fontsize{21}{25}\selectfont #2\par}
  \end{tcolorbox}
  \vspace{2mm}
  \MAIntro{#3}
}

% ---------------------------------------------------------------------------
% Recuadros adicionales sobre el sistema de marco-academic.
% ---------------------------------------------------------------------------
\newtcolorbox{esperado}{ma/base,colback=MASuccessSoft,coltitle=MASuccess,borderline west={2.5pt}{0pt}{MASuccess},title={QUÉ DEBE OBSERVAR Y POR QUÉ}}
\newtcolorbox{alerta}[1][]{ma/base,colback=MADangerSoft,coltitle=MADanger,borderline west={3pt}{0pt}{MADanger},title={SEGURIDAD\ifstrempty{#1}{}{ · #1}}}
\newtcolorbox{diagnostico}[1][]{ma/base,colback=MASurfaceAlt,coltitle=MAInk,borderline west={2.5pt}{0pt}{MAMuted},title={DIAGNÓSTICO\ifstrempty{#1}{}{ · #1}}}
\newtcolorbox{prediccion}[1][]{ma/base,colback=MAWarningSoft,coltitle=MAWarning,borderline west={2.5pt}{0pt}{MAWarning},title={PREDIGA ANTES DE MEDIR\ifstrempty{#1}{}{ · #1}}}
\newtcolorbox{estado}[1][]{ma/base,colback=MATertiarySoft,coltitle=MATertiary,borderline west={2.5pt}{0pt}{MATertiary},title={ESTADO DE VALIDACIÓN\ifstrempty{#1}{}{ · #1}}}

% Paso detallado: acción, dónde, parámetro, por qué, observación, verificación, si falla.
% Un argumento vacío omite su fila (no todos los pasos necesitan las siete partes).
\newtcolorbox{pasobox}[2]{ma/base,breakable=false,colback=MASurface,colframe=MARule,frame hidden,
  borderline={0.4pt}{0pt}{MARule},borderline west={2.5pt}{0pt}{MAPrimary},
  coltitle=MAPrimaryDark,title={PASO #1\enspace·\enspace#2},left=3.5mm,right=3.5mm,top=2.2mm,bottom=2.4mm,
  before skip=6pt,after skip=6pt}
\newcommand{\MApasofila}[2]{\ifstrempty{#2}{}{\textcolor{MAMuted}{\sffamily\bfseries\scriptsize #1} & #2\\}}
\newcommand{\Paso}[7]{%
  \refstepcounter{paso}%
  \begin{pasobox}{\thepaso}{#1}
    \small
    \begin{tabularx}{\linewidth}{@{}>{\raggedright\arraybackslash}p{2.3cm}>{\raggedright\arraybackslash}X@{}}
      \MApasofila{DÓNDE}{#2}%
      \MApasofila{PARÁMETRO}{#3}%
      \MApasofila{POR QUÉ}{#4}%
      \MApasofila{SE OBSERVA}{#5}%
      \MApasofila{VERIFIQUE}{#6}%
      \MApasofila{SI FALLA}{#7}%
    \end{tabularx}
  \end{pasobox}}

% Orden de terminal: comando, propósito, salida esperada, cómo reconocer el fallo, recuperación.
% El comando se escribe tal cual (\detokenize): no use %, # ni llaves desbalanceadas.
\newtcolorbox{comandobox}[1]{ma/base,breakable=false,colback=MASurface,frame hidden,
  borderline={0.4pt}{0pt}{MARule},borderline west={2.5pt}{0pt}{MASecondary},
  left=3.5mm,right=3.5mm,top=2.2mm,bottom=2.4mm,before skip=6pt,after skip=6pt,
  title={\ttfamily\small\color{MAInk}\hyphenpenalty=10000\exhyphenpenalty=10000\$\ #1},
  coltitle=MAInk,fonttitle={},halign title=flush left}
\newcommand{\Comando}[5]{%
  \begin{comandobox}{\detokenize{#1}}
    \small
    \begin{tabularx}{\linewidth}{@{}>{\raggedright\arraybackslash}p{2.3cm}>{\raggedright\arraybackslash}X@{}}
      \MApasofila{PROPÓSITO}{#2}%
      \MApasofila{SE OBSERVA}{#3}%
      \MApasofila{FALLO}{#4}%
      \MApasofila{RECUPERACIÓN}{#5}%
    \end{tabularx}
  \end{comandobox}}

% Rótulos de interfaz y valores: estilo uniforme en todo el manual.
\newcommand{\gui}[1]{\textsf{\textbf{#1}}}
\newcommand{\menu}[1]{\textsf{#1}}
\newcommand{\valor}[1]{\textbf{#1}}
\newcommand{\tecla}[1]{{\setlength{\fboxsep}{1.2pt}\fbox{\sffamily\scriptsize #1}}}
\newcommand{\cmd}[1]{\texttt{\detokenize{#1}}}

% Etiquetas de decisión de preparación de hardware.
\newcommand{\PASS}{{\sffamily\bfseries\color{MASuccess}PASS}}
\newcommand{\COND}{{\sffamily\bfseries\color{MAWarning}CONDICIONAL}}
\newcommand{\BLOQ}{{\sffamily\bfseries\color{MADanger}BLOQUEADO}}

% Casillas para hojas de registro.
\newcommand{\campo}[1][3.2cm]{\rule[-0.4ex]{#1}{0.4pt}}

% ---------------------------------------------------------------------------
% Cuestionarios: una sola fuente (contenido/cuestionarios.tex). El manual del
% estudiante imprime enunciado y opciones; la guía del instructor añade la
% respuesta y su justificación. Activar con \guiainstructortrue.
% ---------------------------------------------------------------------------
\newif\ifguiainstructor
\newcounter{pregunta}
\newcommand{\Pregunta}[7]{% #1 enunciado, #2-#5 opciones A-D, #6 letra correcta, #7 justificación
  \refstepcounter{pregunta}%
  \Needspace{6\baselineskip}%
  \noindent\textbf{\sffamily\small\color{MAPrimaryDark}\thepregunta.}\enspace #1\par\vspace{1pt}
  \begin{enumerate}[label=\textsf{\textbf{\Alph*}}),leftmargin=2.4em,itemsep=0pt,topsep=1pt]
    \item #2
    \item #3
    \item #4
    \item #5
  \end{enumerate}
  \ifguiainstructor
    \vspace{-1pt}{\small\noindent\colorbox{MASuccessSoft}{\sffamily\bfseries\scriptsize\color{MASuccess}\,RESPUESTA #6\,}\enspace #7\par}
  \fi
  \vspace{5pt}}

\setlist[itemize]{leftmargin=1.55em,itemsep=0.25em,topsep=0.35em}
\setlist[enumerate]{leftmargin=1.8em,itemsep=0.35em,topsep=0.35em}

\definecolor{grisclaro}{HTML}{F3F4F6}
\lstset{
	style=ma-code,
	basicstyle=\ttfamily\small,
	backgroundcolor=\color{grisclaro},
	frame=single,
	rulecolor=\color{gray!35},
	breaklines=true,
	columns=fullflexible,
	keepspaces=true,
	showstringspaces=false,
	numbers=none,
	xleftmargin=0.4em,
	xrightmargin=0.4em,
	aboveskip=0.6em,
	belowskip=0.6em
}

% Estilo común de los diagramas TikZ: alto contraste, texto sans legible.
\tikzset{
  bloque/.style={draw=MAPrimaryDark,line width=0.7pt,rounded corners=2pt,fill=MAPrimarySoft,
    align=center,font=\sffamily\scriptsize,minimum height=9mm,inner sep=3pt},
  bloquealt/.style={bloque,draw=MASecondary,fill=MASecondarySoft},
  bloqueinst/.style={bloque,draw=MATertiary,fill=MATertiarySoft},
  bloqueext/.style={bloque,draw=MAMuted,fill=MASurfaceAlt,dashed},
  flecha/.style={-{Stealth[length=2.4mm]},line width=0.8pt,color=MAInk},
  flechadash/.style={flecha,dashed,color=MAMuted},
  rotulo/.style={font=\sffamily\scriptsize,color=MAMuted,align=center},
}
